\documentclass[10pt,a4paper,twocolumn]{article}

% ============================================================
% Single-File LaTeX Editor Test
% Compile with pdfLaTeX.
%
% No external images, fonts, bibliography databases, or data files.
% No labels, cross-references, citations, or hyperlinks.
% ============================================================

% ---------- Page layout ----------
\usepackage[
  a4paper,
  top=18mm,
  bottom=20mm,
  left=17mm,
  right=17mm,
  columnsep=7mm
]{geometry}

% ---------- Basic typography ----------
\usepackage[T1]{fontenc}
\usepackage[utf8]{inputenc}
\usepackage{lmodern}
\usepackage{microtype}

% ---------- Mathematics ----------
\usepackage{amsmath}
\usepackage{amssymb}
\usepackage{amsthm}
\usepackage{mathtools}
\usepackage{bm}

% ---------- Tables, colors, and lists ----------
\usepackage{booktabs}
\usepackage{array}
\usepackage{multirow}
\usepackage{xcolor}
\usepackage{enumitem}

% ---------- Captions and source code ----------
\usepackage[
  font=small,
  labelfont=bf
]{caption}
\usepackage{listings}

% ---------- Header and footer ----------
\usepackage{fancyhdr}

% ---------- Document settings ----------
\setlength{\parindent}{1em}
\setlength{\parskip}{0pt}
\setlength{\columnsep}{7mm}
\setlist{nosep}

% ---------- Colors ----------
\definecolor{codegray}{gray}{0.96}
\definecolor{darkgreen}{rgb}{0.0,0.42,0.12}
\definecolor{darkred}{rgb}{0.55,0.0,0.0}
\definecolor{darkblue}{rgb}{0.0,0.15,0.55}

% ---------- Source-code style ----------
\lstset{
  basicstyle=\ttfamily\scriptsize,
  backgroundcolor=\color{codegray},
  frame=single,
  rulecolor=\color{black!35},
  numbers=left,
  numberstyle=\tiny\color{black!55},
  keywordstyle=\bfseries\color{darkblue},
  commentstyle=\itshape\color{darkgreen},
  stringstyle=\color{darkred},
  breaklines=true,
  showstringspaces=false,
  columns=fullflexible,
  tabsize=2
}

% ---------- Theorem environments ----------
\newtheorem{theorem}{Theorem}
\newtheorem{lemma}[theorem]{Lemma}

\theoremstyle{definition}
\newtheorem{definition}[theorem]{Definition}
\newtheorem{example}[theorem]{Example}

\theoremstyle{remark}
\newtheorem*{remark}{Remark}

% ---------- Custom mathematical commands ----------
\newcommand{\R}{\mathbb{R}}
\newcommand{\E}{\mathbb{E}}
\newcommand{\norm}[1]{\left\lVert #1 \right\rVert}
\newcommand{\abs}[1]{\left\lvert #1 \right\rvert}
\newcommand{\vect}[1]{\bm{#1}}

\DeclareMathOperator{\tr}{tr}
\DeclareMathOperator{\rank}{rank}
\DeclareMathOperator*{\argmin}{arg\,min}

% ---------- Running header ----------
\pagestyle{fancy}
\fancyhf{}
\fancyhead[L]{LaTeX Editor Feature Test}
\fancyhead[R]{Single-File Paper}
\fancyfoot[C]{\thepage}
\renewcommand{\headrulewidth}{0.4pt}

% ---------- Title ----------
\title{
  \textbf{A Single-File Test Document for a \LaTeX{} Editor}\\
  \large Demonstration of Common Paper-Style Typesetting Features
}

\author{
  Example Author\thanks{
    This author note tests footnotes and escaped symbols:
    \%, \$, \&, \#, \_, \{, and \}.
  }\\
  Department of Document Engineering\\
  \texttt{author@example.com}
}

\date{\today}

\begin{document}

% ============================================================
% PAGE 1
% ============================================================

\maketitle

\begin{abstract}
This self-contained document is intended for testing a \LaTeX{}
editor and PDF preview system. It demonstrates text formatting,
sectioning, mathematical notation, theorem environments, lists,
tables, internally generated graphics, source-code listings,
footnotes, colors, headers, footers, and captions. It does not
require external images, external fonts, bibliography databases,
data files, labels, citations, hyperlinks, or cross-references.
Explicit page breaks divide the document into three pages.
\end{abstract}

\noindent
\textbf{Keywords:}
\LaTeX, editor testing, mathematical typesetting, paper layout,
single-file document

\section{Introduction}

Scientific papers commonly combine structured prose, equations,
formal statements, tables, figures, and code fragments. A useful
editor test should therefore exercise more than plain paragraph
rendering.

This paragraph demonstrates several font styles:
\textbf{bold text},
\textit{italic text},
\emph{emphasized text},
\texttt{monospaced text},
\textsc{small capitals},
and {\large larger text}.
Normal text resumes after each local formatting command.

Quotation marks should appear correctly as
``double quotation marks'' and `single quotation marks'.
Dashes can be compared as a hyphen (-), an en dash (--), and an
em dash (---). An ellipsis may be written as \ldots

Common built-in symbols include
\LaTeX{},
\TeX{},
\S,
\copyright,
\pounds,
\dag,
and
\ddag.
Reserved characters are printed by escaping them:
$begin:math:display$
  \\\# \\qquad
  \\\$ \\qquad
  \\\% \\qquad
  \\\& \\qquad
  \\\_ \\qquad
  \\\{ \\qquad
  \\\}\.
$end:math:display$

This sentence contains a footnote.\footnote{
The footnote contains \textit{formatted text}, mathematical notation
such as $begin:math:text$x\^2\+y\^2\=z\^2$end:math:text$, and special characters such as \% and \&.
}

\subsection{Inline Mathematics}

Inline mathematics can be placed inside prose. For example,
$begin:math:text$f\(x\)\=x\^2\+2x\+1$end:math:text$,
$begin:math:text$\\E\[X\]\=\\mu$end:math:text$,
and
$begin:math:text$\\norm\{\\vect\{x\}\}\_2\=\\sqrt\{x\_1\^2\+\\cdots\+x\_n\^2\}$end:math:text$.

Greek letters and common mathematical symbols include
$begin:math:display$
  \\alpha\,\\ \\beta\,\\ \\gamma\,\\ \\Gamma\,\\ \\Delta\,\\ \\theta\,\\ \\lambda\,
  \\mu\,\\ \\pi\,\\ \\rho\,\\ \\sigma\,\\ \\phi\,\\ \\omega\,\\ \\Omega\.
$end:math:display$

Relations and logical symbols include
$begin:math:display$
  a\\leq b\,\\qquad
  b\\geq a\,\\qquad
  x\\neq y\,\\qquad
  p\\approx q\,\\qquad
  A\\subseteq B\,\\qquad
  x\\in\\R\.
$end:math:display$

\subsection{Displayed Equations}

A numbered equation tests fractions, summations, vectors, subscripts,
superscripts, and automatically sized delimiters:

\begin{equation}
  \widehat{\vect{\theta}}
  =
  \argmin_{\vect{\theta}\in\R^d}
  \left\{
    \frac{1}{n}
    \sum_{i=1}^{n}
    \left(
      y_i-\vect{x}_i^{\mathsf T}\vect{\theta}
    \right)^2
    +
    \lambda\norm{\vect{\theta}}_2^2
  \right\}.
\end{equation}

An unnumbered equation can be produced with square brackets:

$begin:math:display$
  e\^\{i\\pi\}\+1\=0\.
$end:math:display$

The following expression demonstrates roots, absolute values,
overlines, hats, and accents:

$begin:math:display$
  \\sqrt\[n\]\{x\}\,
  \\qquad
  \\abs\{x\-y\}\,
  \\qquad
  \\overline\{z\}\,
  \\qquad
  \\widehat\{\\theta\}\,
  \\qquad
  \\widetilde\{f\}\,
  \\qquad
  \\dot\{x\}\,
  \\qquad
  \\ddot\{x\}\.
$end:math:display$

\subsection{Lists}

An itemized list tests bullet rendering and indentation:

\begin{itemize}
  \item plain text and punctuation;
  \item inline mathematics, $begin:math:text$a\^2\+b\^2\=c\^2$end:math:text$;
  \item formatted text such as \textbf{bold} and \textit{italic};
  \item a nested list:
    \begin{itemize}
      \item first nested item;
      \item second nested item.
    \end{itemize}
\end{itemize}

An enumerated list tests custom counters:

\begin{enumerate}[label=(\roman*)]
  \item define the problem;
  \item construct a mathematical model;
  \item calculate the result;
  \item interpret the output.
\end{enumerate}

A description list tests aligned labels:

\begin{description}[style=nextline]
  \item[Input] A finite collection of observations.
  \item[Output] A fitted parameter vector.
  \item[Constraint] No external resource is loaded.
\end{description}

\newpage

% ============================================================
% PAGE 2
% ============================================================

\section{Mathematical Environments}

This page focuses on multiline equations, matrices, piecewise
functions, theorem environments, proofs, and an internally generated
figure.

\subsection{Aligned Equations}

The \texttt{align} environment aligns equations at ampersand
characters:

\begin{align}
  S_n
    &= \sum_{k=1}^{n} k \\
    &= 1+2+\cdots+n \\
    &= \frac{n(n+1)}{2}.
\end{align}

Selected equation numbers can be suppressed with
\verb|\notag|:

\begin{align}
  (a+b)^3
    &= (a+b)(a+b)^2 \notag\\
    &= (a+b)(a^2+2ab+b^2) \notag\\
    &= a^3+3a^2b+3ab^2+b^3.
\end{align}

The \texttt{aligned} environment can be placed inside another
mathematical environment:

$begin:math:display$
  \\begin\{aligned\}
    x\+y \&\= 7\,\\\\
    2x\-y \&\= 5\,\\\\
    3x \&\= 12\,\\\\
    x \&\= 4\,\\qquad y\=3\.
  \\end\{aligned\}
$end:math:display$

\subsection{Piecewise Functions and Delimiters}

A piecewise function is created with the \texttt{cases} environment:

\begin{equation}
  f(x)=
  \begin{cases}
    x^2,       & x\geq 0,\\
    -x,        & -1<x<0,\\
    1,         & x\leq -1.
  \end{cases}
\end{equation}

Different delimiter styles can be tested as follows:

$begin:math:display$
  \(x\)\,\\quad
  \[x\]\,\\quad
  \\\{x\\\}\,\\quad
  \\langle x\\rangle\,\\quad
  \\lfloor x\\rfloor\,\\quad
  \\lceil x\\rceil\.
$end:math:display$

Automatically resized delimiters appear in

$begin:math:display$
  \\left\(
    \\frac\{1\+\\frac\{1\}\{x\}\}\{1\-\\frac\{1\}\{x\}\}
  \\right\)\,
  \\qquad
  \\left\[
    \\sum\_\{k\=1\}\^\{n\} k\^2
  \\right\]\,
  \\qquad
  \\left\\\{
    x\\in\\R \: x\^2\<4
  \\right\\\}\.
$end:math:display$

\subsection{Matrices and Arrays}

Several matrix environments are shown below:

$begin:math:display$
  A\=
  \\begin\{bmatrix\}
    1 \& 2 \& 0\\\\
    0 \& 1 \& 3\\\\
    2 \& 0 \& 1
  \\end\{bmatrix\}\,
  \\qquad
  \\vect\{x\}\=
  \\begin\{pmatrix\}
    x\_1\\\\
    x\_2\\\\
    x\_3
  \\end\{pmatrix\}\.
$end:math:display$

A determinant can be written using vertical delimiters:

$begin:math:display$
  \\det\(A\)
  \=
  \\begin\{vmatrix\}
    1 \& 2 \& 0\\\\
    0 \& 1 \& 3\\\\
    2 \& 0 \& 1
  \\end\{vmatrix\}
  \=13\.
$end:math:display$

An augmented matrix tests vertical rules inside an array:

$begin:math:display$
  \\left\[
  \\begin\{array\}\{ccc\|c\}
    1 \& 1 \& 0 \& 3\\\\
    0 \& 1 \& 1 \& 5\\\\
    2 \& 0 \& 1 \& 4
  \\end\{array\}
  \\right\]\.
$end:math:display$

Matrix operators and transposition can be displayed as

$begin:math:display$
  A\^\{\-1\}A\=I\,
  \\qquad
  \\tr\(A\)\=\\sum\_i a\_\{ii\}\,
  \\qquad
  \\rank\(A\)\\leq\\min\(m\,n\)\,
  \\qquad
  A\^\{\\mathsf T\}\.
$end:math:display$

\subsection{Calculus and Probability}

A definite integral and an infinite series are

\begin{equation}
  \int_{-\infty}^{\infty} e^{-x^2}\,dx=\sqrt{\pi},
  \qquad
  \sum_{k=0}^{\infty}\frac{x^k}{k!}=e^x.
\end{equation}

Derivatives may be written in several forms:

$begin:math:display$
  f\'\(x\)\,
  \\qquad
  \\frac\{df\}\{dx\}\,
  \\qquad
  \\frac\{\\partial f\}\{\\partial x\}\,
  \\qquad
  \\nabla f\(\\vect\{x\}\)\,
  \\qquad
  \\nabla\^2 f\(\\vect\{x\}\)\.
$end:math:display$

Probability notation can be tested with

$begin:math:display$
  \\Pr\(X\\leq x\)\,
  \\qquad
  \\E\[X\]\,
  \\qquad
  \\operatorname\{Var\}\(X\)\,
  \\qquad
  X\\sim\\mathcal\{N\}\(\\mu\,\\sigma\^2\)\.
$end:math:display$

\subsection{Definitions, Theorems, and Proofs}

\begin{definition}[Convex function]
A function $begin:math:text$f\:\\R\^d\\rightarrow\\R$end:math:text$ is convex if, for all
$begin:math:text$\\vect\{x\}\,\\vect\{y\}\\in\\R\^d$end:math:text$ and $begin:math:text$t\\in\[0\,1\]$end:math:text$,

$begin:math:display$
  f\\bigl\(t\\vect\{x\}\+\(1\-t\)\\vect\{y\}\\bigr\)
  \\leq
  t f\(\\vect\{x\}\)\+\(1\-t\)f\(\\vect\{y\}\)\.
$end:math:display$
\end{definition}

\begin{lemma}
For every real number $begin:math:text$x$end:math:text$, the inequality $begin:math:text$x\^2\\geq0$end:math:text$ holds.
\end{lemma}

\begin{proof}
When $begin:math:text$x\\geq0$end:math:text$, the product $begin:math:text$x\\cdot x$end:math:text$ is nonnegative.
When $begin:math:text$x\<0$end:math:text$, the number $begin:math:text$\-x$end:math:text$ is positive and
$begin:math:text$x\^2\=\(\-x\)\^2$end:math:text$. Therefore $begin:math:text$x\^2\\geq0$end:math:text$ in both cases.
\end{proof}

\begin{theorem}[Arithmetic mean minimization]
Let $begin:math:text$x\_1\,\\ldots\,x\_n\\in\\R$end:math:text$. The function
$begin:math:display$
  Q\(c\)\=\\sum\_\{i\=1\}\^\{n\}\(x\_i\-c\)\^2
$end:math:display$
is minimized at
$begin:math:display$
  \\bar\{x\}\=\\frac\{1\}\{n\}\\sum\_\{i\=1\}\^\{n\}x\_i\.
$end:math:display$
\end{theorem}

\begin{proof}
Write
$begin:math:display$
  x\_i\-c\=\(x\_i\-\\bar\{x\}\)\+\(\\bar\{x\}\-c\)\.
$end:math:display$
After expansion,

\begin{align*}
  Q(c)
    &=
    \sum_{i=1}^{n}(x_i-\bar{x})^2\\
    &\quad+
    2(\bar{x}-c)
    \sum_{i=1}^{n}(x_i-\bar{x})\\
    &\quad+
    n(\bar{x}-c)^2.
\end{align*}

The middle sum equals zero. The final term is nonnegative and becomes
zero precisely when $begin:math:text$c\=\\bar\{x\}$end:math:text$. Thus the arithmetic mean minimizes
the squared error.
\end{proof}

\begin{example}
For the values $begin:math:text$2\,4\,6$end:math:text$, the arithmetic mean is $begin:math:text$4$end:math:text$. The squared
error at $begin:math:text$c\=4$end:math:text$ is
$begin:math:display$
  \(2\-4\)\^2\+\(4\-4\)\^2\+\(6\-4\)\^2\=8\.
$end:math:display$
\end{example}

\begin{remark}
This unnumbered environment tests a distinct theorem style without
creating any cross-reference information.
\end{remark}

\subsection{Internally Generated Figure}

The following figure uses only the standard \texttt{picture}
environment. No image file is loaded.

\begin{figure}[htbp]
  \centering
  \setlength{\unitlength}{1mm}

  \begin{picture}(72,34)
    % Coordinate axes
    \put(8,5){\vector(1,0){58}}
    \put(8,5){\vector(0,1){25}}

    % Axis labels
    \put(65,1){\small $begin:math:text$x$end:math:text$}
    \put(3,29){\small $begin:math:text$y$end:math:text$}
    \put(5,1){\small $begin:math:text$0$end:math:text$}

    % Tick marks
    \multiput(18,4)(10,0){5}{\line(0,1){2}}
    \multiput(7,10)(0,5){4}{\line(1,0){2}}

    % A piecewise linear curve
    \put(8,6){\line(1,0){10}}
    \put(18,6){\line(2,1){10}}
    \put(28,11){\line(2,1){10}}
    \put(38,16){\line(2,1){10}}
    \put(48,21){\line(2,1){10}}

    % Point and annotation
    \put(37.3,15.3){\circle*{1.5}}
    \put(40,13){\scriptsize sample point}
  \end{picture}

  \caption{
    A schematic plot generated entirely by \LaTeX{} commands.
    It tests lines, vectors, repeated objects, labels, and captions.
  }
\end{figure}

\newpage

% ============================================================
% PAGE 3
% ============================================================

\section{Tables, Code, and Additional Features}

This page tests floating tables, merged table cells, source-code
listings, verbatim text, colors, minipages, and typographical
details.

\subsection{Basic Table}

The following table uses professional horizontal rules from the
\texttt{booktabs} package:

\begin{table}[htbp]
  \centering
  \caption{Example performance measurements.}
  \small

  \begin{tabular}{
    >{\raggedright\arraybackslash}p{16mm}
    c
    c
    c
  }
    \toprule
    Method
      & Accuracy
      & Time (s)
      & Rank\\
    \midrule
    Baseline
      & $begin:math:text$0\.812$end:math:text$
      & $begin:math:text$2\.31$end:math:text$
      & 3\\
    Proposed
      & $begin:math:text$\\mathbf\{0\.947\}$end:math:text$
      & $begin:math:text$1\.84$end:math:text$
      & $begin:math:text$\\mathbf\{1\}$end:math:text$\\
    Robust
      & $begin:math:text$0\.921$end:math:text$
      & $begin:math:text$2\.05$end:math:text$
      & 2\\
    \bottomrule
  \end{tabular}
\end{table}

\subsection{Grouped Table}

Merged rows and partial rules are demonstrated below:

\begin{table}[htbp]
  \centering
  \caption{Grouped measurements with merged cells.}
  \small

  \begin{tabular}{llrr}
    \toprule
    Group & Condition & Mean & Std. dev.\\
    \midrule

    \multirow{2}{*}{A}
      & Control   & 12.4 & 1.8\\
      & Treatment & 15.7 & 1.3\\

    \cmidrule(lr){1-4}

    \multirow{2}{*}{B}
      & Control   & 11.9 & 2.1\\
      & Treatment & 16.2 & 1.5\\

    \bottomrule
  \end{tabular}
\end{table}

\subsection{Colored Text and Boxes}

Color rendering can be checked with

\begin{center}
  \textcolor{darkblue}{blue text},
  \qquad
  \textcolor{darkgreen}{green text},
  \qquad
  \textcolor{darkred}{red text}.
\end{center}

Background colors can be tested with

\begin{center}
  \colorbox{yellow!35}{highlighted text}
  \qquad
  \fcolorbox{black}{gray!15}{framed text}.
\end{center}

A simple rule and spacing test appears below:

$begin:math:display$
  \\rule\{20mm\}\{0\.4pt\}
  \\qquad
  \\rule\{10mm\}\{2pt\}
  \\qquad
  \\rule\[\-2mm\]\{8mm\}\{6mm\}\.
$end:math:display$

\subsection{Source-Code Listing}

The following listing tests indentation, line numbers, syntax
highlighting, comments, strings, long lines, and underscores:

\begin{lstlisting}[
  language=Python,
  caption={A compact gradient-descent example.}
]
def gradient_descent(
    initial_value,
    learning_rate=0.1,
    tolerance=1e-8,
    max_iterations=1000
):
    """Minimize f(x) = (x - 3)^2."""

    x = float(initial_value)

    for iteration in range(max_iterations):
        gradient = 2.0 * (x - 3.0)
        candidate = x - learning_rate * gradient

        # Stop if the update is sufficiently small.
        if abs(candidate - x) < tolerance:
            return candidate, iteration + 1

        x = candidate

    return x, max_iterations
\end{lstlisting}

The listing is displayed but not executed. Its purpose is to test
verbatim-like source text inside the generated PDF.

\subsection{Inline Verbatim and Verbatim Block}

The \verb|\verb| command can display reserved characters without
escaping them individually:

\begin{center}
  \verb|# $ % & _ { } \|
\end{center}

A larger block can be displayed using the \texttt{verbatim}
environment:

\begin{verbatim}
This text preserves    repeated spaces.
Special characters: # $ % & _ { } \
Commands such as \section are not executed here.
\end{verbatim}

\subsection{Minipages}

Two independently formatted blocks can be placed side by side:

\begin{center}
  \begin{minipage}[t]{0.46\columnwidth}
    \centering
    \textbf{Block A}

    \smallskip

    Left-side content with centered text and an equation:
    $begin:math:display$
      a\^2\+b\^2\=c\^2\.
    $end:math:display$
  \end{minipage}
  \hfill
  \begin{minipage}[t]{0.46\columnwidth}
    \centering
    \textbf{Block B}

    \smallskip

    Right-side content with a short list:
    \begin{itemize}
      \item first;
      \item second.
    \end{itemize}
  \end{minipage}
\end{center}

\subsection{Text Alignment and Spacing}

\begin{flushleft}
This paragraph is explicitly left aligned.
\end{flushleft}

\begin{center}
This line is centered.
\end{center}

\begin{flushright}
This line is right aligned.
\end{flushright}

Horizontal spacing can be inserted with
small space $begin:math:text$a\\\,b$end:math:text$,
medium space $begin:math:text$a\\\:b$end:math:text$,
large space $begin:math:text$a\\\;b$end:math:text$,
quad space $begin:math:text$a\\quad b$end:math:text$,
and double-quad space $begin:math:text$a\\qquad b$end:math:text$.

Vertical spacing can be controlled using commands such as
\verb|\smallskip|, \verb|\medskip|, and \verb|\bigskip|.

\subsection{Conclusion}

This three-page source tests the following editor and renderer
features:

\begin{enumerate}[label=\arabic*.]
  \item paper-style title, abstract, sections, and columns;
  \item bold, italic, monospaced, small-cap, and escaped text;
  \item inline, displayed, aligned, and numbered mathematics;
  \item matrices, cases, calculus, probability, and custom commands;
  \item definitions, lemmas, theorems, examples, and proofs;
  \item lists, tables, merged cells, captions, and floating objects;
  \item an internally generated figure without an image file;
  \item code listings, verbatim blocks, colors, and minipages;
  \item page headers, page footers, footnotes, and explicit breaks.
\end{enumerate}

Because labels, references, citations, hyperlinks, and bibliography
commands are not used, the document does not depend on information
from a previous compilation pass.

\section*{Acknowledgments}

This test document uses only content and resources declared directly
inside the single \texttt{.tex} source file.

\end{document}